\section{Bases teóricas y antecedentes}

\subsection{Moderación, datos y contexto}

La moderación de una plataforma combina reglas, trabajo humano y sistemas automáticos. La predicción de un modelo no resuelve por sí sola la política, la interpretación ni la responsabilidad de una decisión \cite{gorwa2020moderation}. La calidad de un clasificador de lenguaje abusivo depende de cómo se define el daño, cómo se toma la muestra y quién asigna las etiquetas; los errores de esas etapas pueden pasar al modelo \cite{vidgen2020directions}. Los \emph{data statements} y las \emph{data cards} proponen documentar el origen, el idioma, la recolección, el etiquetado, los usos previstos y los límites de un conjunto de datos \cite{bender2018datastatements,pushkarna2022datacards}.

\subsection{Enfoques de solución y trabajos comparables}

El estado del arte resuelve la detección de lenguaje dañino como una cadena y no como la elección aislada de un algoritmo. Primero delimita corpus, población, taxonomía y procedimiento de etiquetado; después representa el texto mediante rasgos léxicos o modelos contextuales; combina clasificadores cuando sus errores son complementarios; evalúa con métricas sensibles al desbalance y pruebas por fenómeno de daño o seguro; y deriva los casos inciertos a revisión humana \cite{vidgen2020directions,bender2018datastatements,sparckjones1972idf,vaswani2017attention,dietterich2000ensemble,rottger2021hatecheck,mozannar2020defer}. Esta secuencia organiza el diseño del presente trabajo mediante datos locales documentados, tres familias de modelos ---clásicos, Transformers y LLM---, ensembles de sus mejores representantes, test con prevalencia explícita y supervisión humana.

Los corpus previos muestran distintas formas de construir la tarea. Wulczyn \emph{et al.} estudian ataques personales a gran escala \cite{wulczyn2017exmachina}; HateXplain añade el blanco del abuso y explicaciones humanas \cite{mathew2021hatexplain}; y HateCheck muestra que una métrica agregada puede ocultar fallas funcionales \cite{rottger2021hatecheck}. La literatura también advierte que las marcas de identidad o las variedades lingüísticas pueden confundirse con toxicidad \cite{sap2019racialbias}. En YouTube, Bertaglia \emph{et al.} analizan lenguaje abusivo en comentarios en inglés \cite{bertaglia2021youtube}; este estudio usa, en cambio, habla transcrita en español y conserva el intervalo temporal de cada evidencia.

La Tabla~\ref{tab:estado_arte} ubica seis antecedentes empleados para interpretar el resultado final. Sus cifras no forman una clasificación común; pues cambian la unidad de análisis, el idioma, el muestreo, la prevalencia, las categorías y la partición. Ninguno de estos datos o modelos interviene en el entrenamiento del proyecto.

\begin{table*}[t]
\centering
\caption{Trabajos relacionados, ámbito de aplicación y relación con el estudio. Fuente: elaboración propia a partir de las publicaciones citadas en cada fila.}
\label{tab:estado_arte}
\scriptsize
\begin{tabularx}{\textwidth}{>{\bfseries\raggedright\arraybackslash}p{20mm}>{\raggedright\arraybackslash}p{31mm}>{\raggedright\arraybackslash}p{57mm}X}
\toprule
\textbf{Trabajo} & \textbf{País o ámbito} & \textbf{Datos, tarea y resultado publicado} & \textbf{Semejanza y diferencia principal} \\
\midrule
DETOXIS \cite{taule2021detoxis} & España; español & 4\,359 comentarios a noticias sobre inmigración; detección binaria de toxicidad; mejor F1 de 0,6461. & Coincide en idioma y daño textual. Cambian la unidad, el dominio temático y el contrato binario. \\
HatEval \cite{basile2019hateval} & Tarea internacional; Twitter en español e inglés & Odio contra mujeres e inmigrantes, con rasgos de agresividad y blanco; mejor macro-F1 en español de 0,730. & Comparte categorías de racismo y género, pero usa mensajes breves y muestreo dirigido al blanco. \\
OffendES \cite{plaza2021offendes} & España; español & 47\,128 comentarios de YouTube, Instagram y Twitter alrededor de influenciadores jóvenes; macro-F1 binario de 0,7839. & Es el antecedente más próximo por YouTube y español; clasifica comentarios escritos, no habla transcrita ni evidencia temporal. \\
EXIST \cite{rodriguez2021exist} & Tarea internacional; español e inglés & Más de 11\,000 textos de Twitter y Gab; identificación y categorización de sexismo; mejor F1 español de 0,7944. & Se aproxima a género/identidad y emplea supervisión experta; no cubre los otros tres daños del proyecto. \\
HateXplain \cite{mathew2021hatexplain} & Twitter y Gab en inglés; autores de India y Alemania & Cerca de 20 mil publicaciones con tres clases, blanco y racionales humanos; macro-F1 de 0,687 para su modelo contextual. & Comparte la necesidad de ligar decisión y evidencia; usa racionales por token, no intervalos de video ni español peruano. \\
NaijaHate \cite{tonneau2024naijahate} & Nigeria; Twitter & 35\,976 tuits y una muestra representativa; precisión promedio (AP) cercana a 0,34 en la muestra aleatoria frente a 0,83--0,90 en conjuntos enriquecidos. & Comparte la preocupación por dominio local y prevalencia; país, idioma, etiqueta y diseño muestral impiden comparar AP de forma directa. \\
\bottomrule
\end{tabularx}
\end{table*}

NaijaHate evidencia que el desempeño aparente y la carga de revisión pueden cambiar entre muestras representativas y enriquecidas \cite{tonneau2024naijahate}. En este proyecto, el resultado principal corresponde al test con prevalencia natural del corpus. La vista secundaria 4:1 usa las mismas predicciones y solo sirve para contextualizar trabajos con muestras enriquecidas.

\subsection{Taxonomía de daño y adaptación peruana}

Para ordenar el espacio de daño se parte de dos distinciones académicas. Waseem \emph{et al.} separan el abuso dirigido a una persona o entidad del dirigido a un grupo, y el abuso explícito del implícito \cite{waseem2017abuse}. Banko \emph{et al.} recomiendan categorías finas, criterios de inclusión y excepciones; su tipología distingue, entre otros, ataque por identidad, insulto, doxeo, agresión sexual y amenaza de violencia \cite{banko2020taxonomy}. El proyecto combina estas distinciones generales con antecedentes sociales e institucionales pertinentes al contexto peruano.

\subsubsection{Aportes para la adaptación al contexto peruano}

Para racismo y discriminación, Juan Carlos Callirgos examinó la construcción del otro en el racismo peruano, y Gonzalo Portocarrero relacionó racismo y mestizaje con prácticas e imaginarios sociales \cite{callirgos1993racismo,portocarrero2009racismo}. Víctor Vich sintetizó la propuesta de Portocarrero mediante el fundamento invisible, el fantasma del patrón y la utopía del blanqueamiento, tres formas de explicar cómo la jerarquización persiste y cambia en el Perú \cite{vich2018dinamicas}. Virginia Zavala y Roberto Zariquiey estudiaron la segregación discursiva justificada mediante supuestas carencias de educación o cultura, mientras el volumen editado por Virginia Zavala y Michele Back reunió análisis sobre los vínculos entre racismo y lenguaje \cite{zavala2007discurso,zavala2017racismo}. En espacios políticos y digitales, Virginia Zavala y Claudia Almeida mostraron que ``motoso'' y ``terruco'' pueden racializar y silenciar a actores políticos; Roberto Brañez Medina documentó el uso de la ortografía para construir las identidades ``amixer'' y ``no-amixer'' y justificar racismo cultural; y Verónica Salem describió cómo el insulto, la parodia, la ironía y el humor reproducen estereotipos raciales en Facebook \cite{almeida2022motoso,branez2012amixer,salem2016amixer}. Estos aportes orientaron las distinciones entre racismo étnico explícito, racialización lingüística, clasismo racial, discriminación regional y racismo encubierto.

Para asuntos de género, Francisco Rodríguez-Sánchez y coautores aportan el esquema computacional de EXIST para identificar sexismo en redes, y Philine Zeinert, Nanna Inie y Leon Derczynski proponen criterios de etiquetado de misoginia en línea \cite{rodriguez2021exist,zeinert2021misogyny}. Christian Monge-Olivarría, Juan Guerra-Corrales y Fernando Bringas-Castro analizan la feminización como insulto y formas de acoso directo, invasión de privacidad y exclusión en Twitter; su estudio aporta mecanismos generales de violencia digital que complementan las fuentes locales \cite{monge2023violencia}. En el ámbito peruano, Denisse Albornoz y Marieliv Flores caracterizan modalidades, consecuencias y respuestas frente a la violencia de género en línea, y la Defensoría del Pueblo del Perú sistematiza su tipología, impactos y vías institucionales de atención \cite{albornoz2018conocer,defensoria2021violenciaenlinea}. Estas fuentes orientan misoginia/acoso de género y, junto con la tipología general, las distinciones sexuales explícitas, de cosificación y no consentidas.

Para orientación sexual e identidad de género, Bharathi Raja Chakravarthi construyó un corpus multilingüe de comentarios de YouTube y referencias de clasificación para homofobia y transfobia \cite{chakravarthi2024homotrans}. En Lima, Jan Marc Rottenbacher de Rojas relacionó conservadurismo e intolerancia a la ambigüedad con homofobia y prejuicio hacia grupos transgénero en una muestra universitaria \cite{rottenbacher2012homofobia}. Carolina Mirian Lovón-Cueva y Marco Lovón-Cueva identificaron en ciberforos peruanos un repertorio de lexemas lesbofóbicos y los interpretaron como violencia simbólica y actos de lenguaje de odio \cite{lovon2022lesbofobia}. Estas fuentes ofrecen un primer acercamiento al tratamiento del odio y la discriminación en entornos peruanos desde los estudios culturales, la sociolingüística, la psicología social, los estudios de género y la documentación institucional \cite{zavala2017racismo,almeida2022motoso,albornoz2018conocer,defensoria2021violenciaenlinea,rottenbacher2012homofobia,lovon2022lesbofobia}. A su vez, fueron usadas para orientar los criterios de inclusión, exclusión y revisión; la cobertura exacta y los nombres finales de las etiquetas de daño son una adaptación operacional del proyecto.

El odio implícito puede expresarse con lenguaje indirecto o codificado \cite{elsherief2021implicit}; ironía y sarcasmo plantean una tarea específica de interpretación \cite{ilic2018irony}. Estudios peruanos muestran que el insulto, la parodia, la ironía y el humor \emph{amixer} \cite{salem2016amixer}, así como la información contextual \cite{bourgeade2024context}, forman parte de la construcción de estereotipos. Estos antecedentes motivan tres señales de revisión, distintas de las categorías de moderación.

Para contenido sexual, la tipología académica y las fuentes peruanas respaldan separar agresión sexual y difusión íntima no consentida \cite{banko2020taxonomy,albornoz2018conocer,defensoria2021violenciaenlinea}; la política de YouTube aporta la frontera operativa entre contenido explícito, sexualización no consentida y excepciones informativas o documentales \cite{youtube2026sexualpolicy}. Con este sustento se formuló una primera taxonomía operacional, que aún requiere validación experta.

\subsection{Modelos de texto y aprendizaje multietiqueta}

\subsubsection{Modelos clásicos}

La frecuencia de término--frecuencia inversa de documento (TF--IDF) representa cada trozo de texto mediante la frecuencia de sus términos y reduce el peso de los que aparecen en muchos documentos \cite{sparckjones1972idf,salton1988tfidf}. Es rápida y conserva vocabulario visible, pero produce vectores dispersos y no modela el contexto. Sobre esos rasgos, la regresión logística estima una probabilidad mediante una frontera lineal; es eficiente y sus coeficientes son interpretables, aunque no capta relaciones no lineales sin crear nuevos rasgos \cite{cox1958logistic}. La máquina de vectores de soporte (SVM) busca un margen amplio entre clases y suele funcionar bien con textos dispersos de alta dimensión; su salida no es una probabilidad y requiere calibración para interpretar el score \cite{cortes1995svm}. Complement Naive Bayes corrige los conteos con información de las clases complementarias, lo que favorece datos textuales desbalanceados; su bajo costo se obtiene a cambio de una fuerte hipótesis de independencia entre rasgos \cite{rennie2003cnb}.

La descomposición en valores singulares (SVD) proyecta TF--IDF a un espacio latente más pequeño \cite{deerwester1990lsa}. Sobre esa proyección, \emph{gradient boosting} combina árboles que corrigen errores sucesivos y puede modelar interacciones no lineales, pero aumenta el ajuste y pierde parte del detalle léxico \cite{friedman2001gbm}. fastText promedia representaciones de palabras y subpalabras antes de aplicar un clasificador lineal; las subpalabras ayudan ante variantes morfológicas y términos poco frecuentes, pero el promedio conserva poco orden y contexto distante \cite{joulin2017fasttext}.

\subsubsection{Transformers compactos}

Los Transformers usan atención para representar cada token según los demás tokens del fragmento \cite{vaswani2017attention}; el modelo \emph{Bidirectional Encoder Representations from Transformers} (BERT) establece el ajuste bidireccional de un modelo preentrenado para una tarea concreta \cite{devlin2019bert}. Sentence-BERT obtiene un vector de oración que permite comparar y clasificar textos con menor costo que procesar cada par por separado \cite{reimers2019sbert}. MiniLM destila relaciones de autoatención de una red mayor y ofrece un codificador compacto \cite{wang2020minilm}. El checkpoint multilingüe \path{sentence-transformers/paraphrase-multilingual-MiniLM-L12-v2} transfiere ese espacio entre idiomas \cite{reimers2020multilingual,hf2026minilmcard}. Su ventaja es el equilibrio entre contexto y costo; sus límites son la longitud finita de entrada y la pérdida de información al resumir el fragmento en un vector.

E5 aprende representaciones con contraste débilmente supervisado y separa consultas y textos mediante prefijos \cite{wang2022e5}. Su extensión multilingüe respalda \path{intfloat/multilingual-e5-small} y el prefijo \texttt{query:} usado en el proyecto \cite{wang2024e5,hf2026e5card}. El entrenamiento contrastivo favorece similitud semántica con un tamaño moderado, pero el resultado depende del formato de entrada y el vector agrupado también actúa como cuello de botella.

\subsubsection{LLM Qwen3 con adaptación LoRA}

Qwen3 es una familia de modelos de lenguaje autorregresivos con versiones densas y de mezcla de expertos, preentrenadas en varios idiomas \cite{qwen2025qwen3}. El estudio usa el checkpoint denso \path{Qwen/Qwen3-0.6B-Base} por su escala, licencia e integración con Transformers \cite{hf2026qwen06bcard}. Una cabeza multietiqueta convierte la representación final en scores de las categorías. LoRA introduce matrices de bajo rango en capas seleccionadas y permite ajustar una fracción pequeña de los parámetros \cite{hu2022lora}. Esta combinación ofrece más capacidad contextual que los modelos compactos, pero exige más memoria y tiempo, y sus scores todavía requieren calibración. \emph{Qwen3--LoRA} designa este clasificador ajustado, no un sinónimo genérico de LLM.

El diseño experimental integra estas representaciones con las estructuras planas, en cascada y jerárquicas descritas en la metodología.

Un fragmento puede activar varias categorías de moderación, por lo que la tarea es multietiqueta \cite{tsoumakas2007multilabel,zhang2014multilabel}. Las jerarquías pueden representar una puerta general y clases hijas \cite{silla2011hierarchical}; los modelos conscientes de jerarquía motivan explotar dependencias entre etiquetas \cite{zhou2020hiagm}.

\subsection{Selección, calibración y abstención}

Precisión, recobrado y F1 resumen aspectos distintos de una matriz de confusión; una sola métrica es insuficiente para todos los tipos de clasificación \cite{sokolova2009metrics}. En datos desbalanceados, las curvas precisión--recobrado describen la clase positiva de forma más informativa que las curvas de característica operativa del receptor (ROC) \cite{saito2015pr,davis2006pr}. La implementación calcula la precisión promedio como \(\AP=\sum_n(R_n-R_{n-1})P_n\), donde \(P_n\) y \(R_n\) corresponden al umbral \(n\) \cite{sklearn2026averageprecision}. Las salidas se calibran por etiqueta con regresión sigmoide \cite{platt1999probabilistic,niculescu2005probabilities}; la calibración de redes modernas requiere comprobación empírica \cite{guo2017calibration}. Por ello, la regla de selección y los umbrales usan el conjunto de validación para limitar el sesgo de selección \cite{cawley2010selection}, y el protocolo congela la selección, la calibración y los umbrales antes de correr los modelos con los datos de test.

Una clasificación selectiva puede abstenerse en casos dudosos \cite{chow1970reject,geifman2017selective}; por ejemplo, el aprendizaje con derivación extiende esa idea al decidir cuándo entregar el caso a una persona experta \cite{mozannar2020defer}. En la política del proyecto, el recobrado de enrutamiento es la fracción de chunks dañinos enviados a alerta o revisión; la precisión es la fracción dañina entre los enviados; la tasa de revisión es la fracción total enviada; y el valor predictivo negativo es la fracción realmente segura entre los casos que permanecen en la ruta automática. El análisis también cuenta los falsos negativos. Estas definiciones se aplican al enrutamiento, no sustituyen las métricas por categoría ni la decisión del supervisor \cite{sokolova2009metrics}.

\subsection{Ontología y trazabilidad terminológica}

Una ontología explicita conceptos compartidos y sus relaciones \cite{gruber1993ontology}; el lenguaje de ontologías web 2 (OWL~2) formaliza clases, propiedades e individuos \cite{w3c2012owl2}, y el sistema simple de organización del conocimiento (SKOS) asocia términos y definiciones controladas \cite{w3c2009skos}. En el vocabulario del proyecto, el \emph{video fuente} es la unidad pública; la \emph{pista de subtítulos}, su transcripción temporal; el \emph{chunk de evidencia}, un intervalo textual con video y marcas de tiempo; y el \emph{etiquetado}, la referencia asignada al chunk. Una \emph{categoría operativa} es una salida multietiqueta entrenable; un \emph{fenómeno fino} o \emph{señal de revisión} es un subtipo de auditoría, no otra salida. \emph{SEGURO} se deriva cuando el texto no activa un daño y no certifica el video completo. Estas definiciones operativas siguen prácticas de documentación y taxonomía \cite{bender2018datastatements,pushkarna2022datacards,banko2020taxonomy}.

La \emph{predicción} reúne los scores de un modelo versionado; la \emph{alerta} supera un umbral o activa revisión; la \emph{decisión humana} acepta, modifica o rechaza la alerta; y la \emph{retroalimentación} o \emph{registro reentrenable} conserva evidencia, predicción y decisión para una iteración futura. El \emph{artefacto de datos} enlaza chunks, fuente, intervalo y partición; el \emph{artefacto de modelo}, arquitectura, pesos, calibradores, umbrales y linaje. Son definiciones propuestas por el trabajo, no categorías jurídicas, y separan automatización y responsabilidad humana \cite{gorwa2020moderation,mozannar2020defer}. La representación legible por máquina aplica parcialmente los principios de datos localizables, accesibles, interoperables y reutilizables (FAIR) \cite{wilkinson2016fair}. El Apéndice~\ref{app:ontologia} presenta el grafo.
